\documentclass[10pt,a4paper]{amsart}
\usepackage{amsmath,amssymb,amsthm,graphicx}
% Retained minimal preamble for the Stacks Project chapter more-algebra.tex (More on Algebra).
% The chapter file pinned-original-source.tex carries only \input{preamble} on its first line, so no
% preamble text lives inside the chapter itself.  The Stacks Project uses one single shared preamble for
% every chapter; the snapshot stages it identically as sources/stacks-chow/preamble.tex and
% sources/stacks-etale-cohomology/preamble.tex.  The block below is copied mechanically, line for line,
% from sources/stacks-chow/preamble.tex lines 14-18 (the Xy-pic setup) and 154-217 (the native theorem
% declarations together with the mathematical macros).  Nothing was added, removed or redefined; the
% wrapper class and the amsmath/amssymb/amsthm/graphicx base packages are supplied by the assembler.
% Chapter: more-algebra.tex -- More on Algebra.
% For commutative diagrams we use Xy-pic
\usepackage[all]{xy}

% We use 2cell for 2-commutative diagrams.
\xyoption{2cell}

% Theorem environments.
%
\theoremstyle{plain}
\newtheorem{theorem}[subsection]{Theorem}
\newtheorem{proposition}[subsection]{Proposition}
\newtheorem{lemma}[subsection]{Lemma}

\theoremstyle{definition}
\newtheorem{definition}[subsection]{Definition}
\newtheorem{example}[subsection]{Example}
\newtheorem{exercise}[subsection]{Exercise}
\newtheorem{situation}[subsection]{Situation}

\theoremstyle{remark}
\newtheorem{remark}[subsection]{Remark}
\newtheorem{remarks}[subsection]{Remarks}

\numberwithin{equation}{subsection}

% Macros
%
\def\lim{\mathop{\mathrm{lim}}\nolimits}
\def\colim{\mathop{\mathrm{colim}}\nolimits}
\def\Spec{\mathop{\mathrm{Spec}}}
\def\Hom{\mathop{\mathrm{Hom}}\nolimits}
\def\Ext{\mathop{\mathrm{Ext}}\nolimits}
\def\SheafHom{\mathop{\mathcal{H}\!\mathit{om}}\nolimits}
\def\SheafExt{\mathop{\mathcal{E}\!\mathit{xt}}\nolimits}
\def\Sch{\mathit{Sch}}
\def\Mor{\mathop{\mathrm{Mor}}\nolimits}
\def\Ob{\mathop{\mathrm{Ob}}\nolimits}
\def\Sh{\mathop{\mathit{Sh}}\nolimits}
\def\NL{\mathop{N\!L}\nolimits}
\def\CH{\mathop{\mathrm{CH}}\nolimits}
\def\proetale{{pro\text{-}\acute{e}tale}}
\def\etale{{\acute{e}tale}}
\def\QCoh{\mathit{QCoh}}
\def\Ker{\mathop{\mathrm{Ker}}}
\def\Im{\mathop{\mathrm{Im}}}
\def\Coker{\mathop{\mathrm{Coker}}}
\def\Coim{\mathop{\mathrm{Coim}}}

% Boxtimes
%
\DeclareMathSymbol{\boxtimes}{\mathbin}{AMSa}{"02}

%
% Macros for moduli stacks/spaces
%
\def\QCohstack{\mathcal{QC}\!\mathit{oh}}
\def\Cohstack{\mathcal{C}\!\mathit{oh}}
\def\Spacesstack{\mathcal{S}\!\mathit{paces}}
\def\Quotfunctor{\mathrm{Quot}}
\def\Hilbfunctor{\mathrm{Hilb}}
\def\Curvesstack{\mathcal{C}\!\mathit{urves}}
\def\Polarizedstack{\mathcal{P}\!\mathit{olarized}}
\def\Complexesstack{\mathcal{C}\!\mathit{omplexes}}
% \Pic is the operator that assigns to X its picard group, usage \Pic(X)
% \Picardstack_{X/B} denotes the Picard stack of X over B
% \Picardfunctor_{X/B} denotes the Picard functor of X over B
\def\Pic{\mathop{\mathrm{Pic}}\nolimits}
\def\Picardstack{\mathcal{P}\!\mathit{ic}}
\def\Picardfunctor{\mathrm{Pic}}
\def\Deformationcategory{\mathcal{D}\!\mathit{ef}}

% --- appended 2026-10-12 (Chief): the project own definitions of the reference/slogan/history
% display environments, copied mechanically from the shared Stacks preamble lines 8-12
% (\usepackage{verbatim} plus the three \newenvironment lines, which the chapter preamble keeps
% verbatim from the project preamble). Nothing added, removed or redefined: these environments
% carry no mathematical content and the project itself typesets them as comments. Without these
% three lines a counted statement whose own text uses them fails with Environment slogan undefined.
\usepackage{verbatim}
\newenvironment{reference}{\comment}{\endcomment}
\newenvironment{slogan}{\comment}{\endcomment}
\newenvironment{history}{\comment}{\endcomment}
\usepackage{cleveref}
\crefname{originallemma}{original lemma}{original lemmas}
\Crefname{originallemma}{Original lemma}{Original lemmas}
\makeatletter
\newlabel{more-algebra-lemma-another-helper-G-ring}{{07PU}{source}{}{sourcepointer-1}{}}
\newlabel{more-algebra-lemma-another-helper-G-ring@cref}{{[originallemma][0][]07PU}{[1][0][]source}{}{}{}}
\makeatother
\crefname{originallemma}{original lemma}{original lemmas}
\Crefname{originallemma}{Original lemma}{Original lemmas}
\makeatletter
\newlabel{algebra-lemma-completion-finite-extension}{{07N9}{source}{}{sourcepointer-2}{}}
\newlabel{algebra-lemma-completion-finite-extension@cref}{{[originallemma][0][]07N9}{[1][0][]source}{}{}{}}
\makeatother
\crefname{originallemma}{original lemma}{original lemmas}
\Crefname{originallemma}{Original lemma}{Original lemmas}
\makeatletter
\newlabel{lemma-derivation-extends}{{07PE}{source}{}{sourcepointer-3}{}}
\newlabel{lemma-derivation-extends@cref}{{[originallemma][0][]07PE}{[1][0][]source}{}{}{}}
\makeatother
\crefname{originallemma}{original lemma}{original lemmas}
\Crefname{originallemma}{Original lemma}{Original lemmas}
\makeatletter
\newlabel{lemma-formally-smooth}{{07EC}{source}{}{sourcepointer-4}{}}
\newlabel{lemma-formally-smooth@cref}{{[originallemma][0][]07EC}{[1][0][]source}{}{}{}}
\makeatother
\crefname{originallemma}{original lemma}{original lemmas}
\Crefname{originallemma}{Original lemma}{Original lemmas}
\makeatletter
\newlabel{lemma-formally-smooth-completion}{{07ED}{source}{}{sourcepointer-5}{}}
\newlabel{lemma-formally-smooth-completion@cref}{{[originallemma][0][]07ED}{[1][0][]source}{}{}{}}
\makeatother
\crefname{originallemma}{original lemma}{original lemmas}
\Crefname{originallemma}{Original lemma}{Original lemmas}
\makeatletter
\newlabel{proposition-fs-regular}{{07PM}{source}{}{sourcepointer-6}{}}
\newlabel{proposition-fs-regular@cref}{{[originallemma][0][]07PM}{[1][0][]source}{}{}{}}
\makeatother
\crefname{originallemma}{original lemma}{original lemmas}
\Crefname{originallemma}{Original lemma}{Original lemmas}
\makeatletter
\newlabel{lemma-quotient-regular}{{07PF}{source}{}{sourcepointer-7}{}}
\newlabel{lemma-quotient-regular@cref}{{[originallemma][0][]07PF}{[1][0][]source}{}{}{}}
\makeatother
\title{Geometric regularity of a completed localisation after a purely inseparable base change}
\author{The Stacks Project Authors}
\date{stacks-project:more-algebra.tex:local-snapshot-20261012; GNU FDL 1.2 or later}
\begin{document}
\maketitle
\setcounter{subsection}{1}
\noindent\textit{Editorial note.} All delivered mathematical text is the author's own, copied from the pinned Stacks Project chapter file more-algebra.tex; no statement or proof step was rewritten, reordered or removed.\par\medskip
\noindent\textit{Editorial note.} The chapter file pinned-original-source.tex carries only \textbackslash{}input\{preamble\} on its first line, so no preamble text lives inside the chapter. The retained file sources/stacks-more-algebra/retained-preamble.tex is a mechanical, line-for-line copy of the project's single shared preamble as staged in the snapshot (Xy-pic setup plus theorem declarations and mathematical macros). Nothing was added, removed or redefined, and the excerpt serves only as preamble, never as a span coordinate.\par\medskip
\noindent\textit{Editorial note.} Both counted spans start at the \textbackslash{}begin line of their environment, so the statement and the author's proof are delivered whole; surrounding prose of the section is deliberately outside the counted spans and is carried as standard prerequisites instead.\par\medskip
\noindent\textit{Editorial note.} The citation fields (source\_id, source\_version, authors, source\_url, license.name, source\_title) are identical to those of every other card this lane has produced from this chapter, so the batch selector treats them as one source.\par\medskip
\noindent\textit{Allowed standard prerequisites (editorial).}\begin{itemize}
\item Completion of a finite extension of Noetherian complete local rings is a product of completions. Algebra, Lemma \textbackslash{}ref\{algebra-lemma-completion-finite-extension\}, cited at proof lines 12514-12515. Rewrites the tensor product with L as a product of the completions of B[x] at the primes over q.
\item A derivation of a ring extends to a localisation and to a completion. Lemma \textbackslash{}ref\{lemma-derivation-extends\}, cited at proof line 12537. Lets the derivation d/dx act on the completed localisation tensored with K, where it sends x - f to a unit.
\item Formal smoothness of the localisation at a maximal ideal and of its completion. Lemmas \textbackslash{}ref\{lemma-formally-smooth\} and \textbackslash{}ref\{lemma-formally-smooth-completion\}, cited at proof lines 12539-12541. Identify the ring tensored with K as formally smooth over A, hence regular.
\item A formally smooth algebra over a field is regular, and a quotient of a regular ring by a nonzerodivisor is regular. Proposition \textbackslash{}ref\{proposition-fs-regular\} and Lemma \textbackslash{}ref\{lemma-quotient-regular\}, cited at proof lines 12542 and 12544. Give regularity of the quotient by x - f, which is the ring whose geometric regularity is claimed.
\end{itemize}\medskip
\ifdefined\hypertarget\hypertarget{sourcepointer-1}{}\fi
\noindent\textit{Source reference.} Original lemma 07PU: focused-seeds/AGP/sources/stacks-snapshot-20261011/tags.txt, source lines 9809--9809. The lemma's own canonical tag, quoted so that readers can resolve the label inside the delivered text; it points at the counted statement itself and carries no proof content.\par\medskip
\ifdefined\hypertarget\hypertarget{sourcepointer-2}{}\fi
\noindent\textit{Source reference.} Original lemma 07N9: focused-seeds/AGP/sources/stacks-snapshot-20261011/tags.txt, source lines 9754--9754. Navigation pointer to a statement the counted proof refers to but does not contain; nothing load-bearing is delegated to it.\par\medskip
\ifdefined\hypertarget\hypertarget{sourcepointer-3}{}\fi
\noindent\textit{Source reference.} Original lemma 07PE: focused-seeds/AGP/sources/stacks-snapshot-20261011/tags.txt, source lines 9794--9794. Navigation pointer to a statement the counted proof refers to but does not contain; nothing load-bearing is delegated to it.\par\medskip
\ifdefined\hypertarget\hypertarget{sourcepointer-4}{}\fi
\noindent\textit{Source reference.} Original lemma 07EC: focused-seeds/AGP/sources/stacks-snapshot-20261011/tags.txt, source lines 9442--9442. Navigation pointer to a statement the counted proof refers to but does not contain; nothing load-bearing is delegated to it.\par\medskip
\ifdefined\hypertarget\hypertarget{sourcepointer-5}{}\fi
\noindent\textit{Source reference.} Original lemma 07ED: focused-seeds/AGP/sources/stacks-snapshot-20261011/tags.txt, source lines 9443--9443. Navigation pointer to a statement the counted proof refers to but does not contain; nothing load-bearing is delegated to it.\par\medskip
\ifdefined\hypertarget\hypertarget{sourcepointer-6}{}\fi
\noindent\textit{Source reference.} Original lemma 07PM: focused-seeds/AGP/sources/stacks-snapshot-20261011/tags.txt, source lines 9802--9802. Navigation pointer to a statement the counted proof refers to but does not contain; nothing load-bearing is delegated to it.\par\medskip
\ifdefined\hypertarget\hypertarget{sourcepointer-7}{}\fi
\noindent\textit{Source reference.} Original lemma 07PF: focused-seeds/AGP/sources/stacks-snapshot-20261011/tags.txt, source lines 9795--9795. Navigation pointer to a statement the counted proof refers to but does not contain; nothing load-bearing is delegated to it.\par\medskip
\begin{lemma}
\label{lemma-another-helper-G-ring}
Let $p$ be a prime number. Let $A$ be a Noetherian complete local domain
with fraction field $K$ of characteristic $p$. Let $\mathfrak q \subset A[x]$
be a maximal ideal lying over the maximal ideal of $A$ and let
$(0) \not = \mathfrak r \subset \mathfrak q$ be a prime lying over
$(0) \subset A$. Then
$A[x]_\mathfrak q^\wedge \otimes_{A[x]} \kappa(\mathfrak r)$ is geometrically
regular over $\kappa(\mathfrak r)$.
\end{lemma}

\begin{proof}
Note that $K \subset \kappa(\mathfrak r)$ is finite.
Hence, given a finite purely
inseparable extension $L/\kappa(\mathfrak r)$ there exists a finite
extension of Noetherian complete local domains $A \subset B$ such that
$\kappa(\mathfrak r) \otimes_A B$ surjects onto $L$.
Namely, you take $B \subset L$
a finite $A$-subalgebra whose field of fractions is $L$. Denote
$\mathfrak r' \subset B[x]$ the kernel of the map
$B[x] = A[x] \otimes_A B \to \kappa(\mathfrak r) \otimes_A B \to L$
so that $\kappa(\mathfrak r') = L$. Then
$$
A[x]_\mathfrak q^\wedge \otimes_{A[x]} L =
A[x]_\mathfrak q^\wedge \otimes_{A[x]} B[x]
\otimes_{B[x]} \kappa(\mathfrak r') =
\prod B[x]_{\mathfrak q_i}^\wedge \otimes_{B[x]} \kappa(\mathfrak r')
$$
where $\mathfrak q_1, \ldots, \mathfrak q_t$ are the primes of $B[x]$
lying over $\mathfrak q$, see
Algebra, Lemma \ref{algebra-lemma-completion-finite-extension}.
Thus we see that it suffices to prove the rings
$B[x]_{\mathfrak q_i}^\wedge \otimes_{B[x]} \kappa(\mathfrak r')$
are regular. This reduces us to showing that
$A[x]_\mathfrak q^\wedge \otimes_{A[x]} \kappa(\mathfrak r)$
is regular in the special case that $K = \kappa(\mathfrak r)$.

\medskip\noindent
Assume $K = \kappa(\mathfrak r)$. In this case we see that
$\mathfrak r K[x]$ is generated by $x - f$ for some $f \in K$
and
$$
A[x]_\mathfrak q^\wedge \otimes_{A[x]} \kappa(\mathfrak r)
=
(A[x]_\mathfrak q^\wedge \otimes_A K)/(x - f)
$$
The derivation $D = \text{d}/\text{d}x$ of $A[x]$ extends to $K[x]$ and
maps $x - f$ to a unit of $K[x]$. Moreover $D$ extends to
$A[x]_\mathfrak q^\wedge \otimes_A K$ by Lemma \ref{lemma-derivation-extends}.
As $A \to A[x]_\mathfrak q^\wedge$ is formally smooth (see
Lemmas \ref{lemma-formally-smooth} and
\ref{lemma-formally-smooth-completion})
the ring $A[x]_\mathfrak q^\wedge \otimes_A K$ is regular by
Proposition \ref{proposition-fs-regular} (the arguments of the
proof of that proposition simplify significantly in this particular case).
We conclude by Lemma \ref{lemma-quotient-regular}.
\end{proof}

\end{document}
